\documentclass[10pt]{article}
\usepackage[utf8]{inputenc}
\usepackage[T1]{fontenc}
\usepackage{amsmath}
\usepackage{amsfonts}
\usepackage{amssymb}
\usepackage[version=4]{mhchem}
\usepackage{stmaryrd}
\usepackage{bbold}
\usepackage{hyperref}
\hypersetup{colorlinks=true, linkcolor=blue, filecolor=magenta, urlcolor=cyan,}
\urlstyle{same}
\usepackage{caption}

\title{Two New Families of Quadratic APN Functions }

\author{Kangquan Li ${ }^{(1)}$, Yue Zhou, Chunlei Li ${ }^{(1)}$, Member, IEEE, and Longjiang Qu ${ }^{(1)}$}
\date{}


%New command to display footnote whose markers will always be hidden
\let\svthefootnote\thefootnote
\newcommand\blfootnotetext[1]{%
  \let\thefootnote\relax\footnote{#1}%
  \addtocounter{footnote}{-1}%
  \let\thefootnote\svthefootnote%
}

%Overriding the \footnotetext command to hide the marker if its value is `0`
\let\svfootnotetext\footnotetext
\renewcommand\footnotetext[2][?]{%
  \if\relax#1\relax%
    \ifnum\value{footnote}=0\blfootnotetext{#2}\else\svfootnotetext{#2}\fi%
  \else%
    \if?#1\ifnum\value{footnote}=0\blfootnotetext{#2}\else\svfootnotetext{#2}\fi%
    \else\svfootnotetext[#1]{#2}\fi%
  \fi
}

\begin{document}
\maketitle
\captionsetup{singlelinecheck=false}


\begin{abstract}
In this paper, we present two new families of APN functions. The first family is in bivariate form ( $x^{3}+x y^{2}+ \left.y^{3}+x y, x^{5}+x^{4} y+y^{5}+x y+x^{2} y^{2}\right)$ over $\mathbb{F}_{2^{m}}^{2}$. It is obtained by adding certain terms of the form $\sum_{i}\left(a_{i} x^{2^{i}} y^{2^{i}}, b_{i} x^{2^{i}} y^{2^{i}}\right)$ to a family of APN functions recently proposed by Göloğlu. The second family has the form $L(z)^{2^{m}+1}+v z^{2^{m}+1}$ over $\mathbb{F}_{2^{3 m}}$, which generalizes a family of APN functions by Bracken et al. from 2011. By calculating the $\Gamma$-rank of the constructed APN functions over $\mathbb{F}_{2^{8}}$ and $\mathbb{F}_{2^{9}}$, we demonstrate that the two families are CCZ-inequivalent to all known families. In addition, the two new families cover two known sporadic APN instances over $\mathbb{F}_{\mathbf{2}^{8}}$ and $\mathbb{F}_{2^{9}}{ }^{9}$, which were found by Edel and Pott in 2009 and by Beierle and Leander in 2021, respectively.
\end{abstract}

Index Terms-APN function, CCZ-equivalence, adding term, bivariate form.

\section*{I. Introduction}
SUBSTITUTION boxes, known as S-boxes, are crucial nonlinear components in block ciphers. They should satisfy a variety of cryptographic criteria [16], such as having low differential uniformity to resist differential attacks [4]. An $n \times n$ S-box can be seen as a function from the finite field $\mathbb{F}_{2^{n}}$ to itself, which is commonly termed an $(n, n)$-function. The differential uniformity of an ( $n, n$ )-function is defined as follows.

Definition 1 ([29]): Given an ( $n, n$ )-function $f: \mathbb{F}_{2^{n}} \rightarrow \mathbb{F}_{2^{n}}$, its differential uniformity is given by\\
$\delta_{f}=\max _{a \in \mathbb{F}_{2^{n}}^{*}, b \in \mathbb{F}_{2^{n}}} \#\left\{z: z \in \mathbb{F}_{2^{n}} \mid f(z+a)+f(z)=b\right\}$.\\
Furthermore, when $\delta_{f}=2$, the function $f$ is called almost perfect nonlinear (APN for short).

APN functions provide the best resistance against differential attacks and also find applications in sequence design and coding theory [17]. In the last three decades, one of the most

\footnotetext{Manuscript received May 18, 2021; revised January 31, 2022; accepted February 22, 2022. Date of publication March 8, 2022; date of current version June 15, 2022. The work of Yue Zhou was supported in part by the Natural Science Foundation of Hunan Province under Grant 2019JJ30030 and in part by the National Natural Science Foundation of China under Grant 62172427. The work of Chunlei Li was supported in part by the Research Council of Norway under Grant 247742/O70 and Grant 311646/O70 and in part by the National Natural Science Foundation of China under Grant 62171479. The work of Longjiang Qu was supported by the National Natural Science Foundation of China under Grant 62032009. (Corresponding author: Longjiang Qu.)

Kangquan Li, Yue Zhou, and Longjiang Qu are with the College of Liberal Arts and Sciences, National University of Defense Technology, Changsha 410073, China, and also with the Hunan Engineering Research Center of Commercial Cryptography Theory and Technology Innovation, Changsha 410013, China (e-mail: \href{mailto:likangquan11@nudt.edu.cn}{likangquan11@nudt.edu.cn}; \href{mailto:gabelozhou@gmail.com}{gabelozhou@gmail.com}; \href{mailto:ljqu_happy@hotmail.com}{ljqu\_happy@hotmail.com}).

Chunlei Li is with the Department of Informatics, University of Bergen, 5020 Bergen, Norway (e-mail: \href{mailto:chunlei.li@uib.no}{chunlei.li@uib.no}).

Communicated by S. Mesnager, Associate Editor At Large for Sequences.\\
Digital Object Identifier 10.1109/TIT.2022.3157810
}\begin{table}[h]
\begin{center}
\captionsetup{labelformat=empty}
\caption{TABLE I\\
All Known APN Monomials Over $\mathbb{F}_{2} n$}
\begin{tabular}{cccc}
\hline
Family & Function & Conditions & Ref. \\
\hline
Gold & $z^{2^{i}+1}$ & $\operatorname{gcd}(i, n)=1$ & $[24]$ \\
Kasami & $z^{2^{2 i}}-2^{i}+1$ & $\operatorname{gcd}(i, n)=1$ & $[26]$ \\
Welch & $z^{2^{t}+3}$ & $n=2 t+1$ & $[20]$ \\
Niho-1 & $z^{2^{t}+2^{t / 2}-1}$ & $n=2 t+1, t$ even & $[19]$ \\
Niho-2 & $z^{2^{t}+2^{(3 t+1) / 2}-1}$ & $n=2 t+1, t$ odd & $[19]$ \\
Inverse & $z^{2^{2 t}-1}$ & $n=2 t+1$ & $[29]$ \\
Dobbertin & $z^{2^{4 i}+2^{3 i}+2^{2 i}+2^{i}-1}$ & $n=5 i$ & $[21]$ \\
\hline
\end{tabular}
\end{center}
\end{table}

important topics in the study of APN functions is to construct new families of APN functions. Several infinite families of APN functions are described in the literature. We summarize all known families of APN functions in Tables I-III according to [9, Table 3], where in Tables I and II the functions are given in univariate form in $z \in \mathbb{F}_{2^{n}}$ and in Table III the functions are given in bivariate form in $(x, y) \in \mathbb{F}_{2}^{2}$. More specifically, Table I lists all known APN monomials and it is conjectured [19] that it is complete, i.e. that there are no other APN monomials (up to equivalence); Table II lists all known families of polynomial APN functions, where $\operatorname{Tr}_{m}^{n}$ denotes the trace function from $\mathbb{F}_{2^{n}}$ to $\mathbb{F}_{2^{m}}$ for any $m \mid n$, i.e., $\operatorname{Tr}_{m}^{n}(z)= z+z^{2^{m}}+\cdots+z^{2^{\left(\frac{n}{m}-1\right) m}}$ and if $m=1, \operatorname{Tr}_{1}^{n}$ becomes the absolutely trace function over $\mathbb{F}_{2^{n}}$, denoted by $\operatorname{Tr}_{n}$ simply; and Table III lists all known families of APN functions proposed originally in bivariate form, where $P_{1}(z)=z^{2^{k}+1}+a z+b$ and $P_{2}(z)=\left(c z^{2^{i}+1}+b z^{2^{i}}+1\right)^{2^{m / 2}+1}+z^{2^{m / 2}+1}$.

In this paper, we propose two new families of quadratic APN functions. The construction of the first family is inspired by Dillon's method [7] and its generalization by Budaghyan and Carlet [10]. Browning et al. in [7] presented a way to construct APN functions of the form


\begin{equation*}
f(z)=z\left(A z^{2}+B^{q}+C z^{2 q}\right)+z^{2}\left(D z^{q}+E z^{2 q}\right)+G z^{3 q} \tag{1}
\end{equation*}


over $\mathbb{F}_{q^{2}}$ with $q=2^{m}$. In particular, Budaghyan and Carlet in [10] obtained a family of APN hexanomials of this form, namely, the family F3 in Table II. Let $f_{1}(z)=A z^{3}+ C z^{2 q+1}+D z^{q+2}+G z^{3 q}$. Recently Li, Li, Helleseth and Qu completely characterized the coefficients for which $f_{1}$ is APN [27]. Later, Chase and Lisoněk [18] proved that when $m \geq 4, f_{1}$ is APN if and only if $f_{1}$ is CCZ-equivalent to one of two instances of Gold functions. This indicates that the APN hexanomials of the form (1) can be seen as the summation of a known APN function $f_{1}(z)$ and $B z^{q+1}+E z^{2(q+1)}$. Such an observation inspired us to seek new APN functions over $\mathbb{F}_{q^{2}}$ by adding special terms of the form $\sum_{i} c_{i} z^{2^{i}(q+1)}, c_{i} \in \mathbb{F}_{q^{2}}^{*}$ to

\begin{table}[h]
\begin{center}
\captionsetup{labelformat=empty}
\caption{TABLE II\\
All Known Polynomial APN Families in Univariate Form Over $\mathbb{F}_{2}{ }^{n}$}
\begin{tabular}{|l|l|l|l|}
\hline
No. & Function & Conditions & Ref. \\
\hline
F1- F2 & $z^{2^{s}+1}+u^{2^{k}-1} z^{2^{i k}+2^{m k+s}}$ & \( \begin{aligned} & n=p k, \operatorname{gcd}(k, 3)=\operatorname{gcd}(s, 3 k)=1, \\ & p \in\{3,4\}, i=s k(\bmod p), m=p-i, \\ & n \geq 12, u \text { primitive in } \mathbb{F}_{2^{n}}^{*} \end{aligned} \) & [11] \\
\hline
F3 & \( \begin{aligned} & s z^{2^{i}(q+1)}+z^{2^{i}+1}+z^{q\left(2^{i}+1\right)} \\ & +c z^{2^{i} q+1}+c^{q} z^{2^{i}+q}+z^{q+1} \end{aligned} \) & \( \begin{aligned} & q=2^{m}, n=2 m, \operatorname{gcd}(i, m)=1 \\ & c \in \mathbb{F}_{2^{n}}, s \in \mathbb{F}_{2^{n}} \backslash \mathbb{F}_{2^{m}}, z^{2^{i}+1}+c z^{2^{i}}+ \\ & c^{q} z+1 \text { has no solution } x \text { with } x^{q+1}=1 \end{aligned} \) & [10] \\
\hline
F4 & $z^{3}+a^{-1} \operatorname{Tr}_{n}\left(a^{3} z^{9}\right)$ & $a \neq 0$ & [12] \\
\hline
F5 & $z^{3}+a^{-1} \operatorname{Tr}_{3}^{n}\left(a^{3} z^{9}+a^{6} z^{18}\right)$ & $3 \mid n, a \neq 0$ & [13] \\
\hline
F6 & $z^{3}+a^{-1} \operatorname{Tr}_{3}^{n}\left(a^{6} z^{18}+a^{12} z^{36}\right)$ & $3 \mid n, a \neq 0$ & [13] \\
\hline
F7-F9 & \( \begin{aligned} & u z^{2^{s}+1}+u^{2^{m}} z^{2^{-m}+2^{m+s}}+ \\ & v z^{2^{-m}+1}+w u^{2^{m}+1} z^{2^{s}+2^{m+s}} \end{aligned} \) & \( \begin{aligned} & n=3 m, \operatorname{gcd}(m, 3)=\operatorname{gcd}(s, 3 m)=1, v, w \in \mathbb{F}_{2^{m}} \\ & v w \neq 1,3 \mid(m+s), u \text { primitive in } \mathbb{F}_{2^{n}}^{*} \end{aligned} \) & [6] \\
\hline
F10 & \( \begin{aligned} & a^{2} z^{2^{2 m+1}+1}+b^{2} z^{2^{m+1}+1}+ \\ & a z^{2^{2 m}+2}+b z^{2^{m}+2}+\left(c^{2}+c\right) z^{3} \end{aligned} \) & \( \begin{aligned} & n=3 m, m \text { odd, } L(z)=a z^{2 m}+b z^{2^{m}}+c z \\ & \text { satisfies the conditions of Lemma } 8 \text { of [8] } \end{aligned} \) & [8] \\
\hline
F11 & \( \begin{aligned} & z^{3}+w z^{2^{i+1}}+w^{2} z^{3 \cdot 2^{m}} \\ & +z^{2^{i+m}}+2^{m} \end{aligned} \) & \( \begin{aligned} & n=2 m, m \text { odd, } 3 \nmid m, w \text { primitive } \\ & \text { in } \mathbb{F}_{2^{2}}, s=m-2,(m-2)^{-1}(\bmod n) \end{aligned} \) & [14] \\
\hline
F12 & $a \operatorname{Tr}_{m}^{n}\left(b z^{3}\right)+a^{q} \operatorname{Tr}_{m}^{n}\left(b^{3} z^{9}\right)$ & \( \begin{aligned} & n=2 m, m \text { odd, } q=2^{m}, a \notin \mathbb{F}_{q}, \\ & b \text { not a cube } \end{aligned} \) & [32] \\
\hline
\end{tabular}
\end{center}
\end{table}

\begin{table}[h]
\begin{center}
\captionsetup{labelformat=empty}
\caption{TABLE III\\
All Known Apn Families in Bivariate Forms Over $\mathbb{F}_{2}^{2}{ }^{m}$}
\begin{tabular}{|l|l|l|l|}
\hline
No. & Function & Conditions & Ref. \\
\hline
F13 & $\left(x y, x^{2^{k}+1}+\alpha x^{\left(2^{k}+1\right) 2^{i}}\right)$ & $\operatorname{gcd}(k, m)=1, m$ even, $\alpha$ non-cubic & [33] \\
\hline
F14 & $\left(x y, x^{2^{3 k}+2^{2 k}}+a x^{2^{2 k}} y^{2^{k}}+b y^{2^{k}+1}\right)$ & $\operatorname{gcd}(k, m)=1, P_{1}$ has no root in $\mathbb{F}_{2^{m}}$ & [30] \\
\hline
F15 & $\left(x y, x^{2^{i}+1}+x^{2^{i+m / 2}} y^{2^{m / 2}}+b x y^{2^{i}}+c y^{2^{i}+1}\right)$ & $m$ even, $\operatorname{gcd}(i, m)=1, P_{2}$ has no root in $\mathbb{F}_{2^{m}}$ & [15] \\
\hline
F16 & $\left(x^{2^{i}+1}+x y^{2^{i}}+y^{2^{i}+1}, x^{2^{2 i}+1}+x^{2^{2 i}} y+y^{2^{2 i}+1}\right)$ & $\operatorname{gcd}(3 i, m)=1$ & [25] \\
\hline
F17 & $\left(x^{2^{i}+1}+x y^{2^{i}}+y^{2^{i}+1}, x^{2^{3 i}} y+x y^{2^{3 i}}\right)$ & $\operatorname{gcd}(3 i, m)=1, m$ odd & [25] \\
\hline
\end{tabular}
\end{center}
\end{table}

known APN families. Note that for the bivariate form over $\mathbb{F}_{q}^{2}$, $\sum_{i} c_{i} z^{2^{i}(q+1)}$ is actually of the form $\sum_{i}\left(a_{i} x^{2^{i}} y^{2^{i}}, b_{i} x^{2^{i}} y^{2^{i}}\right)$, $a_{i}, b_{i} \in \mathbb{F}_{q}^{*}$. It turns out that this approach can give rise to a new family of APN functions over $\mathbb{F}_{2^{m}}^{2}$ by adding the terms $\left(x y, x y+x^{2} y^{2}\right)$ to F16 in Table III.

Theorem 2: Let $m$ be a positive integer with $\operatorname{gcd}(3, m)=1$. Then the function

$$
\begin{aligned}
f(x, y)= & \left(x^{3}+x y^{2}+y^{3}+x y, x^{5}+x^{4} y+\right. \\
& \left.y^{5}+x y+x^{2} y^{2}\right)
\end{aligned}
$$

is APN over $\mathbb{F}_{2}^{2}$.\\[0pt]
The construction of our second family arises from the following APN quadrinomial [6] over $\mathbb{F}_{2^{3 m}}$ :

$$
\begin{aligned}
f(z)= & u z^{2^{s}+1}+u^{2^{m}} z^{2^{-m}+2^{m+s}}+v z^{2^{-m}+1}+ \\
& w u^{2^{m}+1} z^{2^{s}+2^{m+s}}
\end{aligned}
$$

where $\operatorname{gcd}(m, 3)=\operatorname{gcd}(s, 3 m)=1, v, w \in \mathbb{F}_{2^{m}}, v w \neq 1,3 \mid (m+s), u$ is primitive in $\mathbb{F}_{2^{3 m}}$. Assume $w \neq 0$ and take $\gamma=w^{\frac{1}{1-2^{s}}}$, which always exists since $\operatorname{gcd}\left(2^{m}-1,2^{s}-1\right)= 2^{\operatorname{gcd}(m, s)}-1=1$. Then $f(\gamma z)^{2^{m}}=\gamma^{2^{s}+1}\left(L(z)^{2^{m}+1}+\right. \left.(v w+1) z^{2^{m}+1}\right)$, where $L(z)=u^{2^{m}} z^{2^{m+s}}+z$. Note that the equation $L(z)=0$ only has $z=0$ as a solution in $\mathbb{F}_{2^{3 m}}$ since $\operatorname{gcd}\left(2^{3 m}-1,2^{m+s}-1\right)=2^{\operatorname{gcd}(3 m, m+s)}-1=7$ and $u$ is primitive in $\mathbb{F}_{2^{3 m}}$. Thus $L$ is a linearized permutation of $\mathbb{F}_{2^{3 m}}$. Therefore, $f$ is linear equivalent to an APN family of\\
the form $L(z)^{2^{m}+1}+v z^{2^{m}+1}$ with $L$ a permutation binomial and $v \neq 0$. By choosing a linearized permutation trinomial $L$, we propose another new family of APN functions over $\mathbb{F}_{2^{3 m}}$ as follows.

Theorem 3: Let $\operatorname{gcd}(s, m)=1$ and $v \in \mathbb{F}_{2^{m}}^{*}$. Choose $\mu \in \mathbb{F}_{2^{3 m}}^{*}$ such that $\mu^{2^{2 m}+2^{m}+1} \neq 1$ and $L(z)=z^{2^{m+s}}+\mu z^{2^{s}}+z$ permutes $\mathbb{F}_{2^{3 m}}$. Then $f(z)=L(z)^{2^{m}+1}+v z^{2^{m}+1}$ is APN over $\mathbb{F}_{2^{3 m}}$.

Interestingly, the two families in Theorems 2 and 3 both contain previously known sporadic APN instances. According to the code isomorphism test, the family in Theorem 2 includes a sporadic APN function over $\mathbb{F}_{2^{8}}$ (No. 1.9 in Table 9) originally discovered by Edel and Pott in 2009 with the switching method [22]; and the family in Theorem 3 covers a sporadic APN function over $\mathbb{F}_{2^{9}}$, which was recently found by Beierle and Leander in 2021 using a recursive tree search [3]. Here we emphasize that the APN instances in [3] and [22] were not covered by any family before our constructions.

The remainder of this paper is organized as follows. First, in Section II we demonstrate that the proposed two families contain APN functions that are CCZ-inequivalent to those from the known families. Sections III and IV prove the APNness of the families in Theorems 2 and 3, respectively. Section V summarizes the work of this paper and presents some related problems.

\begin{table}[h]
\begin{center}
\captionsetup{labelformat=empty}
\caption{TABLE IV\\
CCZ-Inequivalent Representatives From the Known ApN Families and That in Theorem 2 Over $\mathbb{F}_{2}$ and Their $\Gamma$-Ranks}
\begin{tabular}{|l|l|l|l|}
\hline
No. & Function & $\Gamma$-rank & Ref. \\
\hline
1 & $z^{3}$ & 11818 & Gold \\
\hline
2 & $z^{9}$ & 12370 & Gold \\
\hline
3 & $z^{57}$ & 15358 & Kasami \\
\hline
4 & $z^{3}+z^{17}+u^{48} z^{18}+u^{3} z^{33}+u z^{34}+z^{48}$ & 13200 & F3 \\
\hline
5 & $z^{3}+\operatorname{Tr}_{8}\left(z^{9}\right)$ & 13800 & F4 \\
\hline
6 & $z^{3}+u^{-1} \operatorname{Tr}_{8}\left(u^{3} z^{9}\right)$ & 13842 & F4 \\
\hline
7 & $\left(x y, x^{3}+v y^{12}\right)$ & 13642 & F13 \\
\hline
8 & $\left(x y, x^{12}+x^{4} y^{2}+y^{3}\right)$ & 13700 & F14 \\
\hline
9 & ( $x y, x^{12}+x^{4} y^{2}+v^{7} y^{3}$ ) & 13798 & F14 \\
\hline
10 & $\left(x^{3}+x y^{2}+y^{3}, x^{5}+x^{4} y+y^{5}\right)$ & 13642 & F16 \\
\hline
11 & $\left(x y, x^{3}+x^{2} y+v x^{4} y^{8}+v^{5} y^{3}\right)$ & 13960 & F15 \\
\hline
12 & ( $x^{3}+x y^{2}+y^{3}+x y, x^{5}+x^{4} y+y^{5}+x y+x^{2} y^{2}$ ) & 14034 & Theorem 2 \\
\hline
\end{tabular}
\end{center}
\end{table}

\section*{II. CCZ-Equivalence}
Two functions $f$ and $g$ over $\mathbb{F}_{2^{n}}$ are said to be Carlet-Charpin-Zinoviev ( $C C Z$ ) equivalent [17] if there is an affine permutation of $\mathbb{F}_{2^{n}}^{2}$ that maps the graph $G_{f}= \left\{(z, f(z)): z \in \mathbb{F}_{2^{n}}\right\}$ to the graph $G_{g}=\{(z, g(z)): \left.z \in \mathbb{F}_{2^{n}}\right\}$. CCZ-equivalence is the most general form of equivalence used in the classification of APN functions. To justify that a family of APN functions is indeed new, it is necessary to show that its instances are CCZ-inequivalent to those of the currently known APN families. As a theoretical proof of such a task is rather challenging, in practice one typically chooses to demonstrate that the constructed family contains an APN function that is CCZ-inequivalent to those from another family for small dimensions $n$. Two common approaches are used for this purpose. The first one, known as the code isomorphism test, is to check whether two linear codes associated with the functions are isomorphic. More precisely, for an $(n, n)$-function $f$, its associated code $\mathcal{C}_{f}$ is the linear code defined by the following generating matrix:

$$
\left(\begin{array}{cccc}
1 & 1 & \cdots & 1 \\
0 & u & \cdots & u^{2^{n}}-1 \\
f(0) & f(u) & \cdots & f\left(u^{2^{n}-1}\right)
\end{array}\right),
$$

where $u$ is a primitive element of $\mathbb{F}_{2^{n}}$. It is shown [7], [23] that two functions $f$ and $g$ are CCZ-equivalent if and only if $\mathcal{C}_{f}$ and $\mathcal{C}_{g}$ are isomorphic. The test of code isomorphism is implemented in MAGMA [5], and it can be used to test CCZ-equivalence between two $(n, n)$-functions for small integers $n$. Another common approach is to compare certain CCZ-invariants, i.e., properties that are preserved under CCZ-equivalence, of two $(n, n)$-functions for small integers $n$. If two functions exhibit different values of a given CCZ-invariant, then they must be CCZ-inequivalent. One CCZ-invariant is the $\Gamma$-rank [22]. The $\Gamma$-rank of an $(n, n)$-function $f$ is defined as the rank of the incidence matrix of a design $\operatorname{dev}\left(G_{f}\right)$, whose set of points is $\mathbb{F}_{2^{n}}^{2}$ and whose set of blocks is $\left\{(z+a, f(z)+b): z \in \mathbb{F}_{2^{n}}\right\}$ for $a, b \in \mathbb{F}_{2^{n}}$.

In this section, we demonstrate that our newly constructed APN families are CCZ-inequivalent to all known APN families by comparing the $\Gamma$-ranks of the representatives from all\\
known APN families and ours over $\mathbb{F}_{2^{8}}$ and $\mathbb{F}_{2^{9}}$, see Tables IV and V. ${ }^{1}$ In Table IV, $u$ and $v$ are elements in $\mathbb{F}_{2^{8}}$ and $\mathbb{F}_{2^{4}}$ with minimal polynomials $z^{8}+z^{4}+z^{3}+z^{2}+1$ and $z^{4}+z+1$ over $\mathbb{F}_{2}$, respectively. In Table $\mathrm{V}, u$ denotes an element of $\mathbb{F}_{2^{9}}$ with minimal polynomial $z^{9}+z^{4}+1$ over $\mathbb{F}_{2}$. The last column, denoted by "Ref." for short, refers to the known APN families listed in Tables I, II and III. As shown in Tables IV and V, the two families proposed in this paper contain instances that are not covered by any known family, which indicates that the two families in this paper are new up to CCZ-equivalence. Moreover, according to the code isomorphism test, we observe that the two families also cover two known sporadic APN instances that do not belong to any other known families. As a matter of fact, our first family in Theorem 2 covers one sporadic APN function over $\mathbb{F}_{2^{8}}$ (No. 1.9) in [22, Table 9], which was found by the switching method. The covered APN instance is

$$
\begin{aligned}
& u \operatorname{Tr}_{8}\left(u^{63} z^{3}+u^{252} z^{9}\right)+u^{154} \operatorname{Tr}_{8}\left(u^{68} z^{3}+u^{235} z^{9}\right) \\
& +u^{35} \operatorname{Tr}_{8}\left(u^{216} z^{3}+u^{116} z^{9}\right)+z^{3}
\end{aligned}
$$

where $u$ is defined as in Table IV. In addition, our second APN family in Theorem 3 covers one of the APN functions over $\mathbb{F}_{2^{9}}$ (No. 9) in the dataset [2], which was obtained by a recursive tree search due to Beierle and Leander [3]. The covered APN instance is CCZ-equivalent to

$$
\begin{aligned}
& u^{149} z^{384}+u^{339} z^{320}+u^{498} z^{288}+u^{404} z^{272}+ \\
& u^{125} z^{264}+u^{274} z^{260}+u^{125} z^{258}+u^{432} z^{257}+ \\
& u^{241} z^{192}+u^{14} z^{160}+u^{500} z^{144}+u^{376} z^{136}+ \\
& u^{258} z^{132}+u^{470} z^{130}+u^{430} z^{129}+u^{407} z^{96}+ \\
& u^{317} z^{80}+u^{209} z^{72}+u^{371} z^{68}+u^{77} z^{66}+u^{502} z^{65}+ \\
& u^{464} z^{48}+u^{6} z^{40}+u^{188} z^{36}+u^{498} z^{34}+u^{508} z^{33}+ \\
& u^{199} z^{24}+u^{41} z^{20}+u^{158} z^{18}+u^{218} z^{17}+u^{16} z^{12}+ \\
& u^{7} z^{10}+u^{116} z^{9}+u^{437} z^{6}+u^{502} z^{5}+z^{3}
\end{aligned}
$$

where $u$ is defined as in Table V.

\footnotetext{${ }^{1}$ In Tables IV and V, the $\Gamma$-ranks of the representatives from all known APN families are retrieved from the website \href{https://boolean.h}{https://boolean.h}. \href{http://uib.no/mediawiki/index.php/Tables}{uib.no/mediawiki/index.php/Tables}
}\begin{table}[h]
\begin{center}
\captionsetup{labelformat=empty}
\caption{TABLE V\\
CCZ-Inequivalent Representatives From the Known ApN Families and That in Theorem 3 Over $\mathbb{F}_{2^{9}}$ and Their $\Gamma$-Ranks}
\begin{tabular}{|l|l|l|l|}
\hline
No. & Function & $\Gamma$-rank & Ref. \\
\hline
1 & $z^{3}$ & 38470 & Gold \\
\hline
2 & $z^{5}$ & 41494 & Gold \\
\hline
3 & $z^{17}$ & 38470 & Gold \\
\hline
4 & $z^{13}$ & 58676 & Kasami \\
\hline
5 & $z^{241}$ & 61726 & Kasami \\
\hline
6 & $z^{19}$ & 60894 & Welch \\
\hline
7 & $z^{255}$ & 130816 & Inverse \\
\hline
8 & $z^{3}+\operatorname{Tr}_{9}\left(z^{9}\right)$ & 47890 & F4 \\
\hline
9 & $z^{3}+\operatorname{Tr}_{3}^{9}\left(z^{9}+z^{18}\right)$ & 48428 & F5 \\
\hline
10 & $z^{3}+\operatorname{Tr}_{3}^{9}\left(z^{18}+z^{36}\right)$ & 48460 & F5 \\
\hline
11 & $z^{3}+u^{246} z^{10}+u^{47} z^{17}+u^{181} z^{66}+u^{428} z^{129}$ & 48596 & F10 \\
\hline
12 & $\left(z^{16}+u^{5} z^{2}+z\right)^{9}+u^{73} z^{9}$ & 48558 & Theorem 3 \\
\hline
\end{tabular}
\end{center}
\end{table}

\section*{III. A New Family of ApN Functions Over $\mathbb{F}_{2^{m}}^{2}$}
This section is dedicated to the proof of Theorem 2. We first give several useful lemmas. The following lemma allows us to determine the number of solutions of cubic equations over $\mathbb{F}_{2^{m}}$.

Lemma 1 ([31]): Let $a, b \in \mathbb{F}_{2^{m}}, b \neq 0$ and define

$$
f(z)=z^{3}+a z+b, h(t)=t^{2}+b t+a^{3} .
$$

Let $t_{1}, t_{2}$ be two solutions of $h(t)$ in $\mathbb{F}_{2^{2 m}}$. Then:\\
(1) $f$ has three zeros in $\mathbb{F}_{2^{m}}$ if and only if $\operatorname{Tr}_{m}\left(\frac{a^{3}}{b^{2}}\right)= \operatorname{Tr}_{m}(1), t_{1}$ and $t_{2}$ are cubes in $\mathbb{F}_{2^{m}}$ (resp. $\mathbb{F}_{2^{2 m}}$ ) when $m$ is even (resp. odd).\\
(2) $f$ has exactly one zero in $\mathbb{F}_{2^{m}}$ if and only if $\operatorname{Tr}_{m}\left(\frac{a^{3}}{b^{2}}\right) \neq \operatorname{Tr}_{m}(1)$.\\
(3) $f$ has no zeros in $\mathbb{F}_{2^{m}}$ if and only if $\operatorname{Tr}_{m}\left(\frac{a^{3}}{b^{2}}\right)=\operatorname{Tr}_{m}(1)$, $t_{1}$ and $t_{2}$ are not cubes in $\mathbb{F}_{2^{m}}$ (resp. $\mathbb{F}_{2^{2 m}}$ ) when $m$ is even (resp. odd).\\
The following simple results will be frequently used in the proof of Theorem 2 later.

Lemma 2: Let $\operatorname{gcd}(m, 3)=1$. Then,\\
(1) the equation $z^{3}+z+1=0$ has no solution in $\mathbb{F}_{2^{m}}$;\\
(2) any element $\omega \in \mathbb{F}_{2^{2}} \backslash \mathbb{F}_{2}$ is not cubic in $\mathbb{F}_{2^{m}}$ and $\mathbb{F}_{2^{2 m}}$;\\
(3) the equation $a^{3}+a^{2} b+a+b^{3}+b^{2}+1=0$ holds for $a, b \in \mathbb{F}_{2^{m}}$ if and only if $(a, b)=(1,1)$.\\
Proof: It is clear that any solution of $z^{3}+z+1=0$ is in $\mathbb{F}_{2^{3}}$. Furthermore, any solution in $\mathbb{F}_{2^{m}}$ actually belongs to $\mathbb{F}_{2}$ since $\operatorname{gcd}(m, 3)=1$. Since 0,1 are not solutions of the equation, the first statement follows.

For the second statement, the element $\omega$ is cubic in a finite field $\mathbb{F}_{2^{n}}$ if and only if $\omega^{\frac{2^{n}-1}{3}}=1$, which implies $3 \left\lvert\, \frac{2^{n}-1}{3}\right.$ since $\omega \in \mathbb{F}_{2^{2}} \backslash \mathbb{F}_{2}$. This is equivalent to saying that $\omega \in \mathbb{F}_{2^{2}} \backslash \mathbb{F}_{2}$ is cubic in $\mathbb{F}_{2^{n}}$ if and only if $9 \mid\left(2^{n}-1\right)$, which holds only if $n$ is a multiple of 6 . Since $\operatorname{gcd}(3, m)=1, \omega$ cannot be cubic in $\mathbb{F}_{2^{m}}$ and $\mathbb{F}_{2^{2 m}}$.

Let $B=a^{3}+a^{2} b+a+b^{3}+b^{2}+1$. In the following, we prove that $B=0$ if and only if $(a, b)=(1,1)$. Plugging $a=a_{1}+b$ into $B=0$ and simplifying it, we get


\begin{equation*}
a_{1}^{3}+\left(b^{2}+1\right) a_{1}+(b+1)^{3}=0 . \tag{2}
\end{equation*}


If $b=1$, then $a_{1}=0$ and thus $a=a_{1}+b=1$. If $b \neq$ 1, plugging $a_{1}=(b+1) a_{2}$ into Eq. (2) and simplifying it, we have $a_{2}^{3}+a_{2}+1=0$, which has no solution in $\mathbb{F}_{2^{m}}$ from the first statement of this lemma. Thus $B=0$ if and only if $(a, b)=(1,1)$. \(\square\)

Since the resultant of polynomials will be used in our proof, we now recall some basic facts about it. Given two polynomials $u(x)=a_{m} x^{m}+a_{m-1} x^{m-1}+\cdots+a_{0}$ and $v(x)=b_{n} x^{n}+b_{n-1} x^{n-1}+\cdots+b_{0}$ over a field $K$ with degrees $m$ and $n$, respectively, their resultant $\operatorname{Res}(u, v) \in K$ is the determinant of the following square matrix of order $n+m$ :

$$
\left(\begin{array}{cccccccc}
a_{m} & a_{m-1} & \cdots & a_{0} & 0 & & \cdots & 0 \\
0 & a_{m} & a_{m-1} & \cdots & a_{0} & 0 & \cdots & 0 \\
\vdots & & \ddots & & & & & \vdots \\
0 & \cdots & 0 & a_{m} & a_{m-1} & & \cdots & a_{0} \\
b_{n} & b_{n-1} & \cdots & & b_{0} & 0 & \cdots & 0 \\
0 & b_{n} & b_{n-1} & \cdots & & b_{0} & \cdots & 0 \\
\vdots & & \ddots & & & & \ddots & \vdots \\
0 & \cdots & 0 & b_{n} & b_{n-1} & & \cdots & b_{0}
\end{array}\right) .
$$

For a field $K$ and two polynomials $F(x, y), G(x, y) \in K[x, y]$, we use $\operatorname{Res}_{y}(F, G)$ to denote the resultant of $F$ and $G$ with respect to $y$, which is the resultant of $F$ and $G$ when considered as one polynomial in the single variable $y$. In this case, $\operatorname{Res}_{y}(F, G) \in K[x]$ belongs to the ideal generated by $F$ and $G$. It is known that $F(x, y)=0$ and $G(x, y)=0$ has a common solution $(x, y)$ if and only if $x$ is a solution of $\operatorname{Res}_{y}(F, G)(x)=0$ (see [28, p. 36]).

In the final part of this section, we give the proof of Theorem 2.

Proof of Theorem 2: Since $f$ is a quadratic function with $f(0,0)=(0,0)$, it suffices to show that for any $(a, b) \neq (0,0) \in \mathbb{F}_{2^{m}}^{2}$, the equation


\begin{equation*}
f(x+a, y+b)+f(x, y)+f(a, b)=0 \tag{3}
\end{equation*}


has exactly two solutions $(x, y) \in\{(0,0),(a, b)\}$ in $\mathbb{F}_{2^{m}}^{2}$. By a simple calculation, Eq. (3) is equivalent to the following\\
equation system:

\[
\left\{\begin{array}{l}
F(x, y)=a x^{2}+\left(a^{2}+b^{2}+b\right) x+(a+b) y^{2}+  \tag{4.1}\\
\left(a+b^{2}\right) y=0 \\
G(x, y)=(a+b) x^{4}+b^{2} x^{2}+\left(a^{4}+b\right) x+ \\
b y^{4}+a^{2} y^{2}+\left(a^{4}+a+b^{4}\right) y=0
\end{array}\right.
\]

Firstly, we consider the case $(a, b)=(1,1)$. In this case, Eqs. (4) become

\[
\left\{\begin{array}{l}
x^{2}+x=0  \tag{5.1}\\
x^{2}+y^{4}+y^{2}+y=0
\end{array}\right.
\]

From Eq. (5.1), we know $x \in\{0,1\}$. If $x=0$, plugging it into Eq. (5.2), we have $y^{4}+y^{2}+y=0$ and then $y=$ 0 by Lemma 2 (1). If $x=1$, together with Eq. (5.2), we get $y^{4}+y^{2}+y+1=0$, which means $y=1$. Thus in this case, Eqs. (4) have only $(x, y) \in\{(0,0),(1,1)\}$ as solutions in $\mathbb{F}_{2^{m}}^{2}$.

In the following, we always assume that $(a, b) \neq(1,1)$.\\
With the help of MAGMA, ${ }^{2}$ we obtain the resultant of $F$ and $G$ with respect to $y$ as follows


\begin{gather*}
\operatorname{Res}_{y}(F, G)(x)=\left(a^{3}+a b^{2}+b^{3}\right)^{2} x(x+a) \times \\
H(x, a, b) H(x+a, a, b) \tag{6}
\end{gather*}


where

$$
\begin{aligned}
H(x, a, b)= & x^{3}+\left(a^{2}+a b+a+b^{2}+b+1\right) x+ \\
& a^{3}+a^{2} b+a+b^{3}+b^{2}+1
\end{aligned}
$$

In the sequel we shall show that $\operatorname{Res}_{y}(F, G)(x)=0$ is equivalent to $x(x+a)=0$.

First of all, we have $a^{3}+a b^{2}+b^{3} \neq 0$ for any $(a, b) \neq (0,0) \in \mathbb{F}_{2^{m}}^{2}$. Otherwise, for some element $(a, b) \neq(0,0) \in \mathbb{F}_{2^{m}}^{2}, a^{3}+a b^{2}+b^{3}=0$. If $b=0$, then the above equation becomes $a^{3}=0$, which contradicts the assumption $(a, b) \neq (0,0)$. If $b \neq 0$, then we have $c^{3}+c+1=0$, where $c=\frac{a}{b} \in \mathbb{F}_{2^{m}}$, which contradicts Lemma 2 (1).

In addition, we need to show $H(x, a, b) H(x+a, a, b) \neq 0$. Note that $H(x, a, b)=0$ has the same number of solutions in $\mathbb{F}_{2^{m}}$ as $H(x+a, a, b)=0$. It suffices to show that the equation $H(x, a, b)=0$ has no solution in $\mathbb{F}_{2^{m}}$. Now we consider the equation $H(x, a, b)=0$, i.e.,


\begin{equation*}
x^{3}+A x+B=0, \tag{7}
\end{equation*}


where $A=a^{2}+a b+a+b^{2}+b+1$ and $B=a^{3}+a^{2} b+a+ b^{3}+b^{2}+1$. By Lemma 2 (3), $B \neq 0$ holds under the case $(a, b) \neq(1,1)$. Let $h(t)=t^{2}+B t+A^{3}$. By computation, we have

$$
\frac{A^{3}}{B^{2}}=\frac{C}{B}+\frac{C^{2}}{B^{2}}+1,
$$

where

$$
C=a^{2} b+a^{2}+a b^{2}+a+b^{2}+b .
$$

\footnotetext{${ }^{2}$ we post the MAGMA code to the website \href{https://github.com/yue-zhou-ovgu/APN-functions/blob/main/MAGMA_Proofs.txt}{https://github.com/yue-zhou-ovgu/APN-functions/blob/main/MAGMA\_Proofs.txt}
}Thus $\operatorname{Tr}_{m}\left(\frac{A^{3}}{B^{2}}+1\right)=0$ and the equation $h(t)=0$ has two solutions $t_{1}=C+\omega B$ and $t_{2}=C+\omega^{2} B$ in $\mathbb{F}_{2^{m}}$ (resp. $\mathbb{F}_{2^{2 m}}$ ) if $m$ is even (resp. odd), where $\omega \in \mathbb{F}_{2^{2}} \backslash \mathbb{F}_{2}$. Moreover,

$$
\begin{aligned}
t_{1}= & \omega\left(\omega^{2} C+B\right) \\
= & \omega\left(\omega^{2}\left(a^{2} b+a^{2}+a b^{2}+a+b^{2}+b\right)+\right. \\
& \left.a^{3}+a^{2} b+a+b^{3}+b^{2}+1\right) \\
= & \omega\left(a+\omega b+\omega^{2}\right)^{3}
\end{aligned}
$$

which is not cubic by Lemma 2 (2). Similarly, $t_{2}$ is not cubic. Thus from Lemma 1, the equation $x^{3}+A x+B=0$ has no solution in $\mathbb{F}_{2^{m}}$.

Hence from Eq. (6), we have $x=0$ or $x=a$. Next, we will show that Eqs. (4) have only $(x, y) \in\{(0,0),(a, b)\}$ as solutions.

If $a=0$, then $x=0$. Moreover, Eq. (4.1) and Eq. (4.2) become $b y^{2}+b^{2} y=0$ and $b y^{4}+b^{4} y=0$, respectively. Thus $y=0$ or $y=b$. In the following, we assume that $a \neq 0$.

Case 1: $x=0$. In this case Eqs. (4) become

\[
\left\{\begin{array}{l}
(a+b) y^{2}+\left(a+b^{2}\right) y=0  \tag{8.1}\\
b y^{4}+a^{2} y^{2}+\left(a^{4}+a+b^{4}\right) y=0 .
\end{array}\right.
\]

We now show that Eqs. (8) have only one solution $y=0$ for any $(a, b) \in \mathbb{F}_{2^{m}}^{2} \backslash\{(0,0),(1,1)\}$. If $a=b \neq 1$, then by Eq. (8.1), we get $y=0$. Suppose now that $a \neq b$. If $b=0$, then $a \neq 0$ and Eq. (8.1) becomes $a\left(y^{2}+y\right)=$ 0 , i.e., $y \in\{0,1\}$. In addition, Eq. (8.2) is equivalent to $a^{2} y^{2}+\left(a^{4}+a\right) y=0$. Since $a^{4}+a^{2}+a \neq 0$ for any $a \in \mathbb{F}_{2^{m}}^{*}$, which holds by Lemma 2 (1), $y=1$ is not a solution of Eq. (8.2). Namely, $y=0$ is the unique solution to Eqs. (8). If $a \neq b$ and $b \neq 0$, then from Eq. (8.1), we get $y \in\left\{0, \frac{a+b^{2}}{a+b}\right\}$. Plugging $y=\frac{a+b^{2}}{a+b}$ into Eq. (8.2) and simplifying, we get


\begin{align*}
& a\left(a+b^{2}\right)\left(a^{3}+a b^{2}+b^{3}\right) \times \\
& \left(a^{3}+a^{2} b+a+b^{3}+b^{2}+1\right)=0 . \tag{9}
\end{align*}


Let $c=\frac{a}{b}$. Then $a^{3}+a b^{2}+b^{3}=b^{3}\left(c^{3}+c+1\right) \neq 0$ for any $b \neq 0$ due to Lemma 2 (1). In addition, by Lemma 2 (3), we know that $a^{3}+a^{2} b+a+b^{3}+b^{2}+1 \neq 0$ for any $(a, b) \neq(1,1)$. Thus by Eq. (9), we get $a=b^{2}$, which means that $y=0$ is the unique solution of Eqs. (8).

Case 2: $x=a$. In this case Eqs. (4) become

\[
\left\{\begin{array}{l}
(a+b)(y+b)^{2}+\left(a+b^{2}\right)(y+b)=0  \tag{10.1}\\
b(y+b)^{4}+\left(a^{4}+a+b^{4}\right)(y+b) \\
+a^{2}(y+b)^{2}=0 .
\end{array}\right.
\]

It is clear that $y=b$ is the unique solution of Eqs. (10) by the discussions of the case when $x=0$.

To summarize, Eqs. (4) have only $(x, y) \in\{(0,0),(a, b)\}$ as solutions in $\mathbb{F}_{2^{m}}^{2}$ for any $(a, b) \in \mathbb{F}_{2^{m}}^{2} \backslash\{(0,0)\}$. Therefore, $f$ is APN over $\mathbb{F}_{2^{m}}^{2}$. \(\square\)

\section*{IV. A New Family of ApN Functions Over $\mathbb{F}_{2^{3 m}}$}
In this section, we will show that the univariate function in Theorem 3 is APN. Before that, we give some important lemmas.

Lemma 3: Let $\operatorname{gcd}(m, s)=1, \mu \in \mathbb{F}_{2^{3 m}}^{*}$ satisfy $\mu^{2^{2 m}+2^{m}+1} \neq 1$ and $L_{\beta}(z)=z^{2^{m+s}}+\mu z^{2^{s}}+\beta z$ with $\beta \in \mathbb{F}_{2^{m}}$. Then $L_{0}(z)$ permutes $\mathbb{F}_{2^{3 m}}$. Moreover, $L_{1}$ permutes $\mathbb{F}_{2^{3 m}}$ if and only if $L_{\beta}$ does for any $\beta \in \mathbb{F}_{2^{m}}^{*}$.

Proof: Since $L_{\beta}(x)$ is a linearized polynomial, it suffices to show $L_{\beta}(x)=0$ only has $x=0$ as a solution in $\mathbb{F}_{2^{3 m}}$. Suppose $L_{\beta}(a)=0$ for some $a \in \mathbb{F}_{2^{3 m}}^{*}$. If $\beta=0$, we have $a^{2^{m+s}}+\mu a^{2^{s}}=0$, which contradicts the condition $\mu^{2^{2 m}+2^{m}+1} \neq 1$. If $\beta \neq 0$, let $\epsilon \in \mathbb{F}_{2^{m}}^{*}$ satisfy $\epsilon^{2^{s}-1}=\beta$. Such an element $\epsilon$ always exists since $\operatorname{gcd}\left(2^{m}-1,2^{s}-1\right)= 2^{\operatorname{gcd}(m, s)}-1=1$. Then we have

$$
\begin{aligned}
L_{\beta}(\epsilon x) & =(\epsilon x)^{2^{m+s}}+\mu(\epsilon x)^{2^{s}}+\beta(\epsilon x) \\
& =\epsilon^{2^{s}}\left(x^{2^{m+s}}+\mu x^{2^{s}}+x\right) \\
& =\epsilon^{2^{s}} L_{1}(x)
\end{aligned}
$$

The desired statement thus follows. \(\square\)

Lemma 4: Let $\mu \in \mathbb{F}_{2^{3 m}}^{*}, L(z)=z^{2^{m+s}}+\mu z^{2^{s}}+z$ and $L^{\prime}(z)=z^{2^{m+s}}+\mu^{2^{m}} z^{2^{m}}+z$. Then $L$ permutes $\mathbb{F}_{2^{3 m}}$ if and only if $L^{\prime}$ does.

Proof: For the linear polynomial $L$, it is known that its adjoint polynomial denoted by $L^{*}$ is $L^{*}(z)=z^{2^{2 m-s}}+ \mu^{2^{3 m-s}} z^{2^{3 m-s}}+z$ and $L$ permutes $\mathbb{F}_{2^{3 m}}$ if and only if $L^{*}$ does. Moreover, it is easy to check that $L^{\prime}(z)=\left(L^{*}(z)\right)^{2^{m+s}}$. Thus $L$ permutes $\mathbb{F}_{2^{3 m}}$ if and only if $L^{\prime}$ does. \(\square\)

Let $\operatorname{gcd}(s, m)=1, a \in \mathbb{F}_{2^{3 m}}$ and $v \in \mathbb{F}_{2^{m}}^{*}$. Choose $\mu \in \mathbb{F}_{2^{3 m}}^{*}$ such that $\mu^{2^{2 m}}+2^{m}+1 \neq 1$ and $L(z)=z^{2^{m+s}}+\mu z^{2^{s}}+z$ permutes $\mathbb{F}_{2^{3 m}}$. Define

\[
\left\{\begin{array}{l}
A=L(a) a^{2^{2 m+s}}  \tag{11}\\
B=\left(L(a)^{2^{m}}+\mu^{2^{m}} L(a)\right) a^{2^{m+s}} \\
C=(L(a)+v a) a^{2^{m}} \\
D=\mu L(a)^{2^{m}} a^{2^{s}} \\
E=(L(a)+v a)^{2^{m}} a .
\end{array}\right.
\]

The following lemma is crucial to the main proof in this section.

Lemma 5: Let $A, B, C, D, E$ be defined as in (11), $U_{i}, V_{i}$ with $i=1,2,3,4$ be defined as in (12), shown at the bottom of the next page. Then for any $a \in \mathbb{F}_{2^{3 m}}^{*}$,\\
(i) $A+B+C+D+E=0, A B C D E \neq 0$ and $C+E \neq 0$;\\
(ii) $U_{i} V_{i} \neq 0$ with $i=1,2,3$;\\
(iii) $U_{4}=V_{4}=0$;\\
(iv) $U_{2} V_{1}^{2^{s}}+U_{1} V_{2}^{2^{s}}+U_{3} V_{1}^{2^{s}}+U_{1} V_{3}^{2^{s}}=0$;\\
(v) $U_{2} V_{1}^{2^{s}}+U_{1} V_{2}^{2^{s}} \neq 0$.

Proof: (i) Recall that $L(a)=a^{2^{m+s}}+\mu a^{2^{s}}+a$ and $v \in \mathbb{F}_{2^{m}}$. It is readily seen that

$$
A+B+C+D+E=0 .
$$

Next we show $A B C D E \neq 0$. First, the fact $A D \neq 0$ is clear since $A=L(a) a^{2^{2 m+s}}, D=\mu L(a)^{2^{m}} a^{2^{s}}$ and $L$ permutes $\mathbb{F}_{2^{3 m}}$. Second, let $H(x)=x^{2^{m}}+\mu^{2^{m}} x$. Then $H$ permutes $\mathbb{F}_{2^{3 m}}$ since $\mu^{2^{2 m}}+2^{m}+1 \neq 1$. In addition, it is easy to check that $B=H(L(a)) a^{2^{m+s}}$ and thus $B \neq 0$ for any $a \in \mathbb{F}_{2^{3 m}}^{*}$. Third, $C E \neq 0$ is equivalent to $L(a)+v a=a^{2^{m+s}}+\mu a^{2^{s}}+ (v+1) a \neq 0$, which follows easily from Lemma 3 .

Finally, assume that there exists some $a \in \mathbb{F}_{2^{3 m}}^{*}$ such that $C+E=(L(a)+v a) a^{2^{m}}+(L(a)+v a)^{2^{m}} a=0$. Then\\
$L(a)+v a=\eta a$ for some $\eta \in \mathbb{F}_{2^{m}}^{*}$. Let $\beta=1+v+\eta \in \mathbb{F}_{2^{m}}$. Then $L_{\beta}(a)=0$ for some $a \in \mathbb{F}_{2^{3 m}}^{*}$, which is also impossible by Lemma 3. Thus $C+E \neq 0$ for any $a \in \mathbb{F}_{2^{3 m}}^{*}$.\\
(ii) Let

$$
\begin{aligned}
U= & \mu^{2^{m}} a^{2^{m+s}+1}+a^{2^{2 m+s}+1}+\mu^{2^{m}+1} a^{2^{2 m}+2^{m+s}}+ \\
& a^{2^{m+s}+2^{m}}+\mu^{2^{2 m}+1} a^{2^{2 m+s}+2^{m}}+\mu a^{2^{2 m+s}+2^{2 m}}
\end{aligned}
$$

Plugging the expressions of $A, B, C, D, E$ into those of $U_{1}, U_{2}, U_{3}$ and investigating their factorizations with the help of MAGMA, ${ }^{3}$ we get

\[
\left\{\begin{array}{l}
U_{1}=v a^{2^{m}} C^{2^{2 m}}\left(a^{2^{s}} U\right)^{2^{2 m}}  \tag{13}\\
U_{2}=v a^{2^{m}} C^{2^{2 m}}\left(a^{2^{s}} U\right)^{2^{m}} \\
U_{3}=v a^{2^{m}} C^{2^{2 m}}\left(a^{2^{s}} U\right) .
\end{array}\right.
\]

Moreover, let\\
$V=a^{2^{m}+2^{s}}+\mu^{2^{m}} a^{2^{2 m}+2^{m+s}}+\mu^{2^{2 m}} a^{2^{2 m+s}+2^{m}}+a^{2^{2 m+s}+2^{2 m}}$ and\\
$T=\left(\mu^{2^{2 m}+2^{m}+1}+1\right) a^{2^{s}}+\mu^{2^{2 m}+2^{m}} a+\mu^{2^{2 m}} a^{2^{m}}+a^{2^{2 m}}$.\\
Similarly, plugging the expressions of $A, B, C, D, E$ into those of $V_{1}, V_{2}, V_{3}$ and investigating their factorizations, we obtain

\[
\left\{\begin{array}{l}
V_{1}=v a^{2^{2 m+s+1}+2^{m+s}} L(a) T(a V)^{2^{2 m}}  \tag{14}\\
V_{2}=v a^{2^{2 m+s+1}+2^{m+s}} L(a) T(a V)^{2^{m}} \\
V_{3}=v a^{2^{2 m+s+1}+2^{m+s}} L(a) T(a V) .
\end{array}\right.
\]

Thus in order to show that $U_{i} V_{i} \neq 0$ with $i=1,2,3$, it suffices to prove that $U V T \neq 0$. Let $P=a^{2^{2 m+s}}+\mu^{2^{m}} a^{2^{m+s}}$. Then $P \neq 0$ for any $a \in \mathbb{F}_{2^{3 m}}^{*}$ since the binomial $z^{2^{2 m+s}}+ \mu^{2^{m}} z^{2^{m+s}}$ clearly permutes $\mathbb{F}_{2^{3 m}}$ when $\mu^{2^{2 m}+2^{m}+1} \neq 1$. In addition, we have

$$
V=a^{2^{2 m}} P+a^{2^{m}} P^{2^{m}} .
$$

If there exists some $a \in \mathbb{F}_{2^{3 m}}^{*}$ such that $V=0$, then $P= \tau a^{2^{m}}$ for some $\tau \in \mathbb{F}_{2^{m}}^{*}$. Furthermore, we have $a^{2^{2 m+s}}+ \mu^{2^{m}} a^{2^{m+s}}+\tau a^{2^{m}}=0$, i.e., $a^{2^{m+s}}+\mu a^{2^{s}}+\tau a=0$, which is impossible by Lemma 3. Thus $V \neq 0$. Moreover, $U \neq 0$ due to the crucial observation that

$$
U=V^{2^{2 m}}+\mu V,
$$

which is a permutation in $V$ over $\mathbb{F}_{2^{3 m}}$. Finally, if there exists some $a \in \mathbb{F}_{2^{3 m}}^{*}$ such that $T=0$, then


\begin{align*}
& \left(\mu^{2^{2 m}}+2^{m}+1+1\right) a^{2^{s}}+\mu^{2^{2 m}+2^{m}} a+ \\
& \mu^{2^{2 m}} a^{2^{m}}+a^{2^{2 m}}=0 . \tag{15}
\end{align*}


Raising Eq. (15) to the $2^{m}$-th power, one gets


\begin{align*}
& \left(\mu^{2^{2 m}+2^{m}+1}+1\right) a^{2^{m+s}}+\mu^{2^{2 m}+1} a^{2^{m}}+ \\
& \mu a^{2^{2 m}}+a=0 \tag{16}
\end{align*}


Comparing Eq. (15) and Eq. (16), one can eliminate $\mu^{2^{2 m}+1} a^{2^{m}}+\mu a^{2^{2 m}}$ in Eq. (16) and obtain

$$
\left(\mu^{2^{2 m}+2^{m}+1}+1\right) L(a)=0,
$$

${ }^{3}$ we post the MAGMA code to the website \href{https://github.com/yue-zhou-ovgu/APN-functions/blob/main/MAGMA_Proofs.txt}{https://github.com/yue-zhou-ovgu/APN-functions/blob/main/MAGMA\_Proofs.txt}\\
which is impossible since $\mu^{2^{2 m}+2^{m}+1} \neq 1$ and $L(a) \neq 0$ for any $a \in \mathbb{F}_{2^{3 m}}^{*}$.\\
(iii) It is trivial that $U_{4}=C^{2^{2 m}+2^{m}+1}+E^{2^{2 m}+2^{m}+1}=$ 0 due to the fact $E=C^{2^{m}} a^{1-2^{2 m}}$. The statement $V_{4}=$ 0 holds due to the four equations, shown at the bottom of the page, which can be checked directly.\\
(iv) Plugging Eq. (13) and Eq. (14) into $U_{2} V_{1}^{2^{s}}+U_{1} V_{2}^{2^{s}}+ U_{3} V_{1}^{2^{s}}+U_{1} V_{3}^{2^{s}}$, we get the equation, shown at the bottom of the next page, where

$$
\Delta=v^{2^{s}+1} a^{2^{2 m+2 s+1}+2^{m+2 s}+2^{m}} C^{2^{2 m}} L(a)^{2^{s}} T^{2^{s}}
$$

Moreover, it is easy to check


\begin{equation*}
a^{2^{2 m+s}} U^{2^{2 m}}+a^{2^{m+s}} U^{2^{m}}+a^{2^{s}} U=0 \tag{17}
\end{equation*}


and


\begin{equation*}
a^{2^{2 m}} V^{2^{2 m}}+a^{2^{m}} V^{2^{m}}+a V=0 \tag{18}
\end{equation*}


By multiplying both sides of Eq. (17) with $V^{2^{2 m+s}}$, we get $a^{2^{m+s}} U^{2^{m}} V^{2^{2 m+s}}+a^{2^{s}} U V^{2^{2 m+s}}=a^{2^{2 m+s}} U^{2^{2 m}} V^{2^{2 m+s}}$.

In addition, by raising Eq. (18) to its $2^{s}$-th power and multiplying both sides with $U^{2^{2 m}}$, we obtain\\
$a^{2^{m+s}} U^{2^{2 m}} V^{2^{m+s}}+a^{2^{s}} U^{2^{2 m}} V^{2^{s}}=a^{2^{2 m+s}} U^{2^{2 m}} V^{2^{2 m+s}}$.\\
Together with the above two equations, we have

$$
\begin{aligned}
& a^{2^{m+s}} U^{2^{m}} V^{2^{2 m+s}}+a^{2^{m+s}} U^{2^{2 m}} V^{2^{m+s}} \\
+ & a^{2^{s}} U V^{2^{2 m+s}}+a^{2^{s}} U^{2^{2 m}} V^{2^{s}}=0
\end{aligned}
$$

and therefore

$$
U_{2} V_{1}^{2^{s}}+U_{1} V_{2}^{2^{s}}+U_{3} V_{1}^{2^{s}}+U_{1} V_{3}^{2^{s}}=0
$$

(v) By direct computations, we have

$$
\begin{aligned}
& U_{2} V_{1}^{2^{s}}+U_{1} V_{2}^{2^{s}} \\
= & \Delta\left(a^{2^{m+s}} U^{2^{m}}\left(a^{2^{2 m}} V^{2^{2 m}}\right)^{2^{s}}+a^{2^{2 m+s}} U^{2^{2 m}}\left(a^{2^{m}} V^{2^{m}}\right)^{2^{s}}\right) \\
= & \Delta a^{2^{2 m+s}+2^{m+s}}\left(U V^{2^{m+s}}+U^{2^{m}} V^{2^{s}}\right)^{2^{m}}
\end{aligned}
$$

where $\Delta=v^{2^{s}+1} a^{2^{2 m+2 s+1}+2^{m+2 s}+2^{m}} C^{2^{2 m}} L(a)^{2^{s}} T^{2^{s}}$. Thus in order to prove $U_{2} V_{1}^{2^{s}}+U_{1} V_{2}^{2^{s}} \neq 0$, it suffices to show $U V^{2^{m+s}}+U^{2^{m}} V^{2^{s}} \neq 0$. If there exists some $a \in \mathbb{F}_{2^{3 m}}^{*}$ such that $U V^{2^{m+s}}+U^{2^{m}} V^{2^{s}}=0$, then $U=\gamma V^{2^{s}}$ for some $\gamma \in \mathbb{F}_{2^{m}}^{*}$. In addition, since $U=V^{2^{2 m}}+\mu V$, we obtain

$$
V^{2^{2 m}}+\mu V+\gamma V^{2^{s}}=0
$$

Let $\epsilon \in \mathbb{F}_{2^{m}}^{*}$ satisfy $\epsilon^{1-2^{s}}=\gamma$. Replacing $V$ with $\epsilon V$ in the above equation and multiplying both sides of the equation with $\epsilon^{-1}$, we have $V^{2^{2 m}}+\mu V+V^{2^{s}}=0$. By raising the above equation to its $2^{m}$-th power, we get

$$
V^{2^{m+s}}+\mu^{2^{m}} V^{2^{m}}+V=0
$$

which is impossible by Lemma 4 and the fact that $L$ permutes $\mathbb{F}_{2^{3 m}}$.

Now we give the proof of Theorem 3.\\
Proof of Theorem 3: Since $f$ is a quadratic function with $f(0)=0$, it suffices to show that for any $a \in \mathbb{F}_{2^{3 m}}^{*}$, the equation $f(a z+a)+f(a z)+f(a)=0$ has exactly two solutions $z \in\{0,1\}$ in $\mathbb{F}_{2^{3 m}}$. More specifically, we need to show that the equation


\begin{align*}
& (L(a z)+L(a))^{2^{m}+1}+L(a z)^{2^{m}+1}+L(a)^{2^{m}+1}+ \\
& v(a z+a)^{2^{m}+1}+v(a z)^{2^{m}+1}+v a^{2^{m}+1} \\
= & A z^{2^{2 m+s}}+B z^{2^{m+s}}+C z^{2^{m}}+D z^{2^{s}}+E z \\
= & 0 \tag{19}
\end{align*}



\begin{align*}
& \left\{\begin{aligned}
U_{1} & =D^{2^{2 m}} E^{2^{m}+1}+A C^{2^{2 m}} E^{2^{m}}+B^{2^{m}} C^{2^{2 m}+1} \\
U_{2} & =A^{2^{2 m}} E^{2^{m}+1}+B C^{2^{2 m}} E^{2^{m}}+C^{2^{2 m}+1} D^{2^{m}} \\
U_{3} & =B^{2^{2 m}} E^{2^{m}}+1+C^{2^{2 m}} D E^{2^{m}}+A^{2^{m}} C^{2^{2 m}}+1 \\
U_{4} & =C^{2^{2 m}}+2^{m}+1+E^{2^{2 m}}+2^{m}+1 \\
V_{1} & =A^{2^{2 m}+2} C^{2^{m}}+A B C^{2^{m}} D^{2^{2 m}}+A B^{2^{m}+1} E^{2^{2 m}}+A^{2} D^{2^{m}} E^{2^{2 m}} \\
V_{2} & =A^{2^{2 m}+2} E^{2^{m}}+A B D^{2^{2 m}} E^{2^{m}}+A^{2^{2 m}+1} B^{2^{m}} C+A C D^{2^{2 m}}+2^{m} \\
V_{3} & =A^{2^{2 m}+1} B^{2^{m}} E+A B^{2^{m}+1} C^{2^{2 m}}+A^{2} C^{2^{2 m}} D^{2^{m}}+A D^{2^{2 m}+2^{m}} E \\
V_{4} & =\left(B^{2^{m}+1}+A D^{2^{m}}\right)\left(A B^{2^{2 m}}+D^{2^{2 m}+1}\right)+\left(A^{2^{2 m}}+1+B D^{2^{2 m}}\right)\left(A^{2^{m}+1}+B^{2^{m}} D\right) .
\end{aligned}\right.  \tag{12}\\
& \begin{aligned}
& B^{2^{m}+1}+A D^{2^{m}}=\left(L(a)^{2^{m}}+\mu^{2^{m}} L(a)\right)^{2^{m}+1} a^{2^{2 m+s}+2^{m+s}}+\mu^{2^{m}} L(a)^{2^{2 m}+1} a^{2^{2 m+s}+2^{m+s}} \\
&=a^{2^{2 m+s}+2^{m+s}} L(a)^{2^{m}}\left(L(a)^{2^{2 m}}+\mu^{2^{2 m}} L(a)^{2^{m}}+\mu^{2^{2 m}+2^{m}} L(a)\right) . \\
& A B^{2^{2 m}}+D^{2^{2 m}+1}=L(a) a^{2^{2 m+s}}\left(L(a)^{2^{m}}+\mu^{2^{m}} L(a)\right)^{2^{2 m}} a^{2^{s}}+\mu^{2^{2 m}+1} L(a)^{2^{m}+1} a^{2^{2 m+s}+2^{s}} \\
&=a^{2^{2 m+s}+2^{s}} L(a)\left(\mu L(a)^{2^{2 m}}+\mu^{2^{2 m}+1} L(a)^{2^{m}}+L(a)\right) . \\
& \begin{aligned}
A^{2^{2 m}+1}+B D^{2^{2 m}} & =L(a)^{2^{2 m}+1} a^{2^{2 m+s}+2^{m+s}}+\left(L(a) \mu^{2^{m}}+L(a)^{2^{m}}\right) a^{2^{m+s}} \mu^{2^{2 m}} L(a) a^{2^{2 m+s}} \\
& =a^{2^{2 m+s}+2^{m+s}} L(a)\left(L(a)^{2^{2 m}}+\mu^{2^{2 m}} L(a)^{2^{m}}+\mu^{2^{2 m}+2^{m}} L(a)\right) .
\end{aligned} \\
& A^{2^{m}+1}+B^{2^{m}} D=L(a)^{2^{m}+1} a^{2^{2 m+s}+2^{s}}+\left(L(a) \mu^{2^{m}}+L(a)^{2^{m}}\right)^{2^{m}} a^{2^{2 m+s}} \mu L(a)^{2^{m}} a^{2^{s}} \\
&=a^{2^{2 m+s}+2^{s}} L(a)^{2^{m}}\left(\mu L(a)^{2^{2 m}}+\mu^{2^{2 m}+1} L(a)^{2^{m}}+L(a)\right) .
\end{aligned}
\end{align*}


where $A, B, C, D, E$ are defined as in (11), has exactly two solutions $z \in\{0,1\}$. By Lemma 5 (i), we know $A B C D E \neq$ 0 for any $a \in \mathbb{F}_{2^{3 m}}^{*}$. Raising Eq. (19) to its $2^{m}$-th power and its $2^{2 m}$-th power, we have


\begin{align*}
& A^{2^{m}} z^{2^{s}}+B^{2^{m}} z^{2^{2 m+s}}+C^{2^{m}} z^{2^{2 m}}+ \\
& D^{2^{m}} z^{2^{m+s}}+E^{2^{m}} z^{2^{m}}=0 \tag{20}
\end{align*}


and


\begin{align*}
& A^{2^{2 m}} z^{2^{m+s}}+B^{2^{2 m}} z^{2^{s}}+C^{2^{2 m}} z+ \\
& D^{2^{2 m}} z^{2^{2 m+s}}+E^{2^{2 m}} z^{2^{2 m}}=0 \tag{21}
\end{align*}


respectively. In the following, we will use the method of elimination twice and finally acquire two equations (23) and (26).

After computing the summation of Eq. (19) multiplied by $E^{2^{m}}$ and Eq. (20) multiplied by $C$ and simplifying it, we get


\begin{align*}
& \left(A E^{2^{m}}+B^{2^{m}} C\right) z^{2^{2 m+s}}+C^{2^{m}+1} z^{2^{2 m}}+ \\
& \left(B E^{2^{m}}+C D^{2^{m}}\right) z^{2^{m+s}}+\left(D E^{2^{m}}+A^{2^{m}} C\right) z^{2^{s}}+ \\
& E^{2^{m}+1} z=0 \tag{22}
\end{align*}


By computing the summation of Eq. (21) multiplied by $E^{2^{m}}+1$ and Eq. (22) multiplied by $C^{2^{2 m}}$, we have

$$
U_{1} z^{2^{2 m+s}}+U_{2} z^{2^{m+s}}+U_{3} z^{2^{s}}+U_{4} z^{2^{2 m}}=0
$$

i.e.,


\begin{equation*}
U_{1} z^{2^{2 m+s}}+U_{2} z^{2^{m+s}}+U_{3} z^{2^{s}}=0 \tag{23}
\end{equation*}


where $U_{i}$ with $i=1,2,3,4$ are defined as in (12) and $U_{4}=$ 0 by Lemma 5 (iii).

After computing the summation of Eq. (19) multiplied by $B^{2^{m}}$ and Eq. (20) multiplied by $A$ and simplifying it, we get


\begin{align*}
& A C^{2^{m}} z^{2^{2 m}}+\left(B^{2^{m}+1}+A D^{2^{m}}\right) z^{2^{m+s}}+ \\
& \left(B^{2^{m}} C+A E^{2^{m}}\right) z^{2^{m}}+\left(B^{2^{m}} D+A^{2^{m}+1}\right) z^{2^{s}}+ \\
& B^{2^{m}} E z=0 \tag{24}
\end{align*}


By calculating the summation of Eq. (20) multiplied by $D^{2^{2 m}}$ and Eq. (21) multiplied by $B^{2^{m}}$, we obtain


\begin{align*}
& \left(C^{2^{m}} D^{2^{2 m}}+B^{2^{m}} E^{2^{2 m}}\right) z^{2^{2 m}}+ \\
& \left(D^{2^{2 m}}+2^{m}+A^{2^{2 m}} B^{2^{m}}\right) z^{2^{m+s}}+D^{2^{2 m}} E^{2^{m}} z^{2^{m}}+ \\
& \left(A^{2^{m}} D^{2^{2 m}}+B^{2^{2 m}}+2^{m}\right) z^{2^{s}}+B^{2^{m}} C^{2^{2 m}} z=0 \tag{25}
\end{align*}


Now computing the summation of Eq. (24) multiplied by $\left(D^{2^{2 m}}+2^{m}+A^{2^{2 m}} B^{2^{m}}\right)$ and Eq. (25) multiplied by ( $B^{2^{m}+1}+ \left.A D^{2^{m}}\right)$, we have

$$
V_{1} z^{2^{2 m}}+V_{2} z^{2^{m}}+V_{3} z+V_{4} z^{2^{s}}=0
$$

i.e.,


\begin{equation*}
V_{1}^{2^{s}} z^{2^{2 m+s}}+V_{2}^{2^{s}} z^{2^{m+s}}+V_{3}^{2^{s}} z^{2^{s}}=0 \tag{26}
\end{equation*}


where $V_{i}$ with $i=1,2,3,4$ are defined as in (12) and $V_{4}=$ 0 by Lemma 5 (iii).

Now the summation of Eq. (23) multiplied by $V_{1}^{2^{s}}$ and Eq. (26) multiplied by $U_{1}$ gives\\
$\left(U_{2} V_{1}^{2^{s}}+V_{2}^{2^{s}} U_{1}\right) z^{2^{m+s}}+\left(U_{3} V_{1}^{2^{s}}+V_{3}^{2^{s}} U_{1}\right) z^{2^{s}}=0$,\\
i.e.,

$$
\left(U_{2} V_{1}^{2^{s}}+V_{2}^{2^{s}} U_{1}\right)\left(z^{2^{m+s}}+z^{2^{s}}\right)=0
$$

since $U_{2} V_{1}^{2^{s}}+V_{2}^{2^{s}} U_{1}+U_{3} V_{1}^{2^{s}}+V_{3}^{2^{s}} U_{1}=0$ by Lemma 5 (iv). Moreover, by Lemma 5 (v), $U_{2} V_{1}^{2^{s}}+V_{2}^{2^{s}} U_{1} \neq 0$ and thus $z^{2^{m+s}}+z^{2^{s}}=0$. In other words, $z \in \mathbb{F}_{2^{m}}$. Plugging it into Eq. (19), we get

$$
(C+E)\left(z^{2^{s}}+z\right)=0
$$

By Lemma 5 (i), $C+E \neq 0$ and thus $z^{2^{s}}+z=0$. Then $z \in \mathbb{F}_{2^{\operatorname{gcd}(m, s)}}=\mathbb{F}_{2}$.

In conclusion, Eq. (19) has only two solutions $z \in\{0,1\}$ in $\mathbb{F}_{2^{3 m}}$ for any $a \in \mathbb{F}_{2^{3 m}}^{*}$ and then $f$ is APN over $\mathbb{F}_{2^{3 m}}$.

Remark 1: The existence of the parameter $\mu \in \mathbb{F}_{2^{3 m}}^{*}$ in Theorem 3 for any $m$ is a crucial problem to the significance of the theorem. During the review process, by using techniques from algebraic varieties over finite fields, Bartoli et al. [1] proved that when $m \geq 3$, for the particular case $s=1$, there always exists $\mu \in \mathbb{F}_{2^{3 m}}^{*}$ satisfying $\mu^{2^{2 m}}+2^{m}+1 \neq 1$ such that $L(z)=z^{2^{m+s}}+\mu z^{2^{s}}+z$ permutes $\mathbb{F}_{2^{3 m}}$. That is to say, Theorem 3 indeed always produces APN functions over $\mathbb{F}_{2^{3 m}}$ for any positive integer $m \geq 3$.

\section*{V. Conclusion and Further Work}
In this paper, we obtained two new infinite families of APN functions over $\mathbb{F}_{2^{2 m}}$ and $\mathbb{F}_{2^{3 m}}$. It is demonstrated that our APN families are CCZ-inequivalent to all known infinite families of APN functions by comparing their $\Gamma$-ranks over $\mathbb{F}_{2^{8}}$ or $\mathbb{F}_{2^{9}}$. Furthermore, it is worth mentioning that our two APN families cover two known sporadic APN instances over $\mathbb{F}_{2^{8}}$ and $\mathbb{F}_{2^{9}}$, which were found by Edel and Pott [22] in 2009 using the switching method and Beierle and Leander [3] in 2021 using a recursive tree search, respectively.

Note that the newly found APN family over $\mathbb{F}_{2^{m}}^{2}$ can be expressed using

$$
\begin{gathered}
f(z)=z^{3}+A z^{3 q}+B z^{2 q+1}+C z^{q+2}+D z^{5}+E z^{5 q}+ \\
F z^{4 q+1}+G z^{q+4}+H z^{q+1}+I z^{2(q+1)}
\end{gathered}
$$

over $\mathbb{F}_{q^{2}}$ with $q=2^{m}$. Thus the next question is whether it is possible to obtain new infinite families of APN functions over $\mathbb{F}_{q^{2}}$ of the above form, or more generally,

$$
\begin{aligned}
f(z)= & z\left(A z^{2}+B z^{4}+C z^{q}+D z^{2 q}+E z^{4 q}\right)+ \\
& z^{2}\left(G z^{4}+H z^{q}+I z^{2 q}+J z^{4 q}\right)+
\end{aligned}
$$

$$
\begin{aligned}
& U_{2} V_{1}^{2^{s}}+U_{1} V_{2}^{2^{s}}+U_{3} V_{1}^{2^{s}}+U_{1} V_{3}^{2^{s}} \\
= & \Delta\left(a^{2^{m+s}} U^{2^{m}}\left(a^{2^{2 m}} V^{2^{2 m}}\right)^{2^{s}}+a^{2^{2 m+s}} U^{2^{2 m}}\left(a^{2^{m}} V^{2^{m}}\right)^{2^{s}}+a^{2^{s}} U\left(a^{2^{2 m}} V^{2^{2 m}}\right)^{2^{s}}+a^{2^{2 m+s}} U^{2^{2 m}}(a V)^{2^{s}}\right) \\
= & a^{2^{2 m+s}} \Delta\left(a^{2^{m+s}} U^{2^{m}} V^{2^{2 m+s}}+a^{2^{m+s}} U^{2^{2 m}} V^{2^{m+s}}+a^{2^{s}} U V^{2^{2 m+s}}+a^{2^{s}} U^{2^{2 m}} V^{2^{s}}\right) .
\end{aligned}
$$

$$
\begin{aligned}
& z^{4}\left(K z^{q}+L z^{2 q}+M z^{4 q}\right)+ \\
& z^{q}\left(N z^{2 q}+P z^{4 q}\right)+Q z^{6 q}
\end{aligned}
$$

which was discussed in [10]. In addition, for the newly found APN family over $\mathbb{F}_{2^{3 m}}$, it is important to study the CCZ-equivalence among the APN functions $f(z)= \left(z^{2^{m+s}}+\mu z^{2^{s}}+z\right)^{2^{m}+1}+v z^{2^{m}+1}$ with different parameters $\mu, s, v$ and to determine a lower bound on the number of CCZ-inequivalent APN functions over $\mathbb{F}_{2^{3 m}}$ of this form.

\section*{Acknowledgment}
The authors would like to thank deeply the Associate Editor Prof. Sihem Mesnager and the anonymous referees for their valuable comments and suggestions that highly improved both the technical and the editorial qualities of this article.

\section*{References}
[1] D. Bartoli, M. Calderini, O. Polverino, and F. Zullo, "On the infiniteness of a family of APN functions," J. Algebra, vol. 598, pp. 68-84, May 2022.\\[0pt]
[2] C. Beierle and G. Leander, "New instances of quadratic APN functions in small dimension," Zenodo, Dec. 2021, doi: 10.5281/zenodo. 4738942.\\[0pt]
[3] C. Beierle and G. Leander, "New instances of quadratic APN functions," IEEE Trans. Inf. Theory, vol. 68, no. 1, pp. 670-678, Jan. 2022.\\[0pt]
[4] E. Biham and A. Shamir, "Differential cryptanalysis of DES-like cryptosystems," J. Cryptol., vol. 4, no. 1, pp. 3-72, Jan. 1990.\\[0pt]
[5] W. Bosma, C. Cannon, and C. Playoust, "The MAGMA algebra system I: The user language," J. Symbolic Comput., vol. 24, nos. 1-9, pp. 235-265, Sep. 1997.\\[0pt]
[6] C. Bracken, E. Byrne, N. Markin, and G. McGuire, "A few more quadratic APN functions," Cryptogr. Commun., vol. 3, pp. 43-53, Mar. 2011.\\[0pt]
[7] K. A. Browning et al., "APN polynomials and related codes," J. Combinatorics, Inf. Syst. Sci., vol. 34, nos. 1-4, pp. 135-159, 2009.\\[0pt]
[8] L. Budaghyan, M. Calderini, C. Carlet, R. S. Coulter, and I. Villa, "Constructing APN functions through isotopic shifts," IEEE Trans. Inf. Theory, vol. 66, no. 8, pp. 5299-5309, Aug. 2020.\\[0pt]
[9] L. Budaghyan, M. Calderini, and I. Villa, "On equivalence between known families of quadratic APN functions," Finite Fields Appl., vol. 66, Sep. 2020, Art. no. 101704.\\[0pt]
[10] L. Budaghyan and C. Carlet, "Classes of quadratic APN trinomials and hexanomials and related structures," IEEE Trans. Inf. Theory, vol. 54, no. 5, pp. 2354-2357, May 2008.\\[0pt]
[11] L. Budaghyan, C. Carlet, and G. Leander, "Two classes of quadratic APN binomials inequivalent to power functions," IEEE Trans. Inf. Theory, vol. 54, no. 9, pp. 4218-4229, Sep. 2008.\\[0pt]
[12] L. Budaghyan, C. Carlet, and G. Leander, "Constructing new APN functions from known ones," Finite Fields Appl., vol. 15, pp. 150-159, Apr. 2009.\\[0pt]
[13] L. Budaghyan, C. Carlet, and G. Leander, "On a construction of quadratic APN functions," in Proc. IEEE Inf. Theory Workshop, Oct. 2009, pp. 374-378.\\[0pt]
[14] L. Budaghyan, T. Helleseth, and N. Kaleyski, "A new family of APN quadrinomials," IEEE Trans. Inf. Theory, vol. 66, no. 11, pp. 7081-7087, Nov. 2020.\\[0pt]
[15] M. Calderini, L. Budaghyan, and C. Carlet, "On known constructions of APN and AB functions and their relation to each other," Math. Znanosti, vol. 25, no. 60, pp. 79-105, 2021.\\[0pt]
[16] C. Carlet, Boolean Functions for Cryptography Coding Theory. Cambridge, U.K.: Cambridge Univ. Press, 2021.\\[0pt]
[17] C. Carlet, P. Charpin, and V. Zinoviev, "Codes, bent functions and permutations suitable for DES-like cryptosystems," Des., Codes, Cryptogr., vol. 15, no. 2, pp. 125-156, 1998.\\[0pt]
[18] B. Chase and P. Lisonák, "Kim-type APN functions are affine equivalent to gold functions," Cryptography Commun., vol. 13, no. 6, pp. 981-993, Nov. 2021.\\[0pt]
[19] H. Dobbertin, "Almost perfect nonlinear power functions on $\operatorname{GF}\left(2^{n}\right)$ : The Niho case," Inf. Comput., vol. 151, nos. 1-2, pp. 57-72, May 1999.\\[0pt]
[20] H. Dobbertin, "Almost perfect nonlinear power functions on GF ( $2^{n}$ ): The Welch case," IEEE Trans. Inf. Theory, vol. 45, no. 4, pp. 1271-1275, May 1999.\\[0pt]
[21] H. Dobbertin, "Almost perfect nonlinear power functions on $\operatorname{GF}\left(2^{n}\right)$ : A new case for $n$ divisible by 5," in Finite Fields Application. Berlin, Germany: Springer, 2001, pp. 113-121.\\[0pt]
[22] Y. Edel and A. Pott, "A new almost perfect nonlinear function which is not quadratic," Adv. Math. Commun., vol. 3, no. 1, pp. 59-81, Feb. 2009.\\[0pt]
[23] Y. Edel and A. Pott, "On the equivalence of nonlinear functions," in Enhancing Cryptography Primitives With Technology From Error Correcting Codes. Amsterdam, The Netherlands: IOS Press, 2009, pp. 87-103.\\[0pt]
[24] R. Gold, "Maximal recursive sequences with 3 -valued recursive crosscorrelation functions (Corresp.)," IEEE Trans. Inf. Theory, vol. IT-14, no. 1, pp. 154-156, Jan. 1968.\\[0pt]
[25] F. Göloǧlu, "Biprojective almost perfect nonlinear functions," IEEE Trans. Inf. Theory, 2022, doi: 10.1109/TIT.2022.3157798.\\[0pt]
[26] T. Kasami, "The weight enumerators for several classes of subcodes of the 2nd order binary Reed-Müller codes," Inf. Control, vol. 18 pp. 369-394, May 1971.\\[0pt]
[27] K. Li, C. Li, T. Helleseth, and L. Qu, "A complete characterization of the APN property of a class of quadrinomials," IEEE Trans. Inf. Theory, vol. 67, no. 11, pp. 7535-7549, Nov. 2021.\\[0pt]
[28] R. Lidl and H. Niederreiter, Finite Fields. Cambridge, U.K.: Cambridge Univ. Press, 1997.\\[0pt]
[29] K. Nyberg, "Differentially uniform mappings for cryptography," in Advance Cryptology. Springer, 1993, pp. 55-64.\\[0pt]
[30] H. Taniguchi, "On some quadratic APN functions," Des., Codes Cryptogr., vol. 87, no. 9, pp. 1973-1983, Sep. 2019.\\[0pt]
[31] K. S. Williams, "Note on cubics over $\operatorname{G} F\left(2^{n}\right)$ and $\operatorname{G} F\left(3^{n}\right)$," J. Number Theory, vol. 7, no. 4, pp. 361-365, 1975.\\[0pt]
[32] L. Zheng, H. Kan, Y. Li, J. Peng, and D. Tang, "Constructing new APN functions through relative trace functions," 2021, arXiv:2101.11535.\\[0pt]
[33] Y. Zhou and A. Pott, "A new family of semifields with 2 parameters," Adv. Math., vol. 234, pp. 43-60, Feb. 2013.

Kangquan Li received the B.S., M.S., and Ph.D. degrees from the National University of Defense Technology, Changsha, China, in 2015, 2017, and 2021, respectively. Currently, he is a Lecturer with the National University of Defense Technology. His research interests include cryptography and coding theory.

Yue Zhou was born in Taiyuan, China, in 1984. He received the B.S. and M.S. degrees in mathematics from the National University of Defense Technology, Changsha, China, in 2006 and 2009, respectively, and the Ph.D. degree in mathematics from Otto-von-Guericke University Magdeburg, Germany, in 2013. Currently, he is a Professor with the National University of Defense Technology. His research interests include finite geometries, coding theory, and finite fields.

Chunlei Li (Member, IEEE) received the Ph.D. degree from the University of Bergen, Norway, in 2014. He was a Post-Doctoral Researcher at the University of Stavanger, Norway, from 2015 to 2017, and a Researcher at the University of Bergen from 2017 to 2018. Since 2018, he has been an Associate Professor with the Department of Informatics, University of Bergen. His research interests include sequence design, coding theory, and cryptography. He was the Program Co-Chair of the workshops Mathematical Methods for Cryptography in 2017 and Sequences and Their Applications (SETA) in 2020. He served as a Program Committee Member for several workshops, including WAIFI18, SETA18, BFA20/21/22, IWSDA19/22, and WCC22.

Longjiang Qu received the B.A. and Ph.D. degrees in mathematics from the National University of Defense Technology, Changsha, China, in 2002 and 2007, respectively. He is currently a Professor with the College of Liberal Arts and Sciences, National University of Defense Technology. His research interests include cryptography and coding theory.


\end{document}